\sectionwithtoclink{Fourier Pricing}

When dealing with more complex models for asset prices than the Black--Scholes model, the fact that there is no known analytical
density for the underlying at time $T$ becomes a common problem, as for a european derivative the calculation $\mathbb{E}[w(S_T)]$ is
no longer a straightforward integration exercise. For a subset of these more complicated models we've discovered analytical
expressions for their characteristic functions. Using these as the means to price derivatives is exactly the task we will address over the next section.

\subsection{The Lewis Formula}
The following result is due to Lewis~\cite{lewis2001}.
Let $X_t$ be a process with the ``flipped'' characteristic function $u\to\psi(-u)$, assume that this function is analytic and regular on the strip $\mathcal{S}_{-\psi}$.
Let $w(x)\geq0$ with generalised Fourier transform $\hat{w}$ on strip of regularity $\mathcal{S}_w$. For
$\alpha\in \{v\in\mathbb{R} : \{x+i v:x\in\mathbb{R}\}\subseteq \mathcal{S}_M=\mathcal{S}_{-\psi} \cap \mathcal{S}_w\}=A$
\begin{align}
\mathbb{E}[w(X_T)] = \frac{1}{2\pi}\int_{-\infty+i\alpha}^{\infty+i\alpha}\psi(-z)\hat{w}(z) dz
\end{align}
\begin{proof}
Assume that $A$ is non-empty.
Let $\alpha\in A$.
By the inversion theorem and the choice of $\alpha$:
\begin{align}
\mathbb{E}[w(X_T)]&=\int_{-\infty}^{\infty}w(x)\frac{1}{2\pi}\int_{-\infty+i\alpha}^{\infty+i\alpha}\exp(-izx)\psi(-z) dzdx
\end{align}
We check that the product is integrable via Tonelli's theorem:
\begin{equation}\int_{-\infty}^{\infty} \int_{-\infty}^{\infty}
|w(x)| \cdot e^{-\alpha x} \cdot |\psi(-(u+i\alpha))|\, du\, dx < \infty\end{equation}
since $w(x)\exp(-\alpha x)\in L^1$ and $\psi(-(u+i\alpha))\in L^1$
since $\alpha\in A$, the integral is finite.
Thus we may apply Fubini's theorem:
\begin{align}
(3)&=\int_{-\infty+i\alpha}^{\infty+i\alpha}\psi(-z)\frac{1}{2\pi}\int_{-\infty}^{\infty}\exp(-izx)w(x) dxdz \\
&=\frac{1}{2\pi}\int_{-\infty+i\alpha}^{\infty+i\alpha}\psi(-z)\hat{w}(z) dz\quad
\end{align}
\end{proof}
Now, turning this into something useful for finance, we can pick a model of our choice with a nice characteristic function, a positive payoff function, and then calculate this integral.

\subsection{Pricing Options via the Covered Call/Cash Secured Put}
Throughout this section we work at time $t=0$ with maturity $T$, pricing today's option price.
When one wants to apply this method to pricing options directly, it turns out that the available set $\mathcal{S}_M$ will narrow and widen depending
on one's choice of parameters, which will make the whole task of calibrating a volatility surface based on a model more difficult.
To circumvent this problem we will not price a call directly but rather price a covered-call.\\
A covered call is what one gets by buying a forward on the underlying and selling a call on it, ``covering'' one's downside, giving the following payoff at time $T$:
\begin{align}
e^{X_T}-\max(0,e^{X_T}-K)=\min(e^{X_T},K) \\
Z = e^{-rT}\mathbb{E}^{\mathbb{Q}}[\min(e^{X_T},K)]
\end{align}
then creating the call by replication.
\begin{equation}C = e^{-rT}F_T-Z
\end{equation}
noting that by put-call parity
\begin{equation}P = C-e^{-rT}F +e^{-rT}K = e^{-rT}K-Z
\end{equation}
so let's find the strip of regularity for this payoff
\begin{align}
    z= \beta+i \alpha \in S_f \Leftrightarrow \int_{-\infty}^{\infty}e^{-\alpha x}|\min(e^{x},K)|dx<\infty\\
\int_{-\infty}^{\infty}e^{-\alpha x}|\min(e^{x},K)|dx = \int_{-\infty}^{\log(K)}e^{(1-\alpha) x} dx+\int_{\log(K)}^{\infty}e^{-\alpha x}K dx
\end{align}
Both integrals are positive, we simply need to figure out when they are finite, we do this rather simply with monotone convergence:

\textbf{First integral:}\\
For $\alpha =1$ we clearly diverge.\\
For $\alpha \neq 1$:
\begin{equation}\int_{-\infty}^{\log(K)}e^{(1-\alpha) x} dx= \lim_{n\to \infty}[-\frac{e^{(1-\alpha) x}}{1-\alpha}]^{\log(K)}_{-n}<\infty \Leftrightarrow \alpha<1
\end{equation}

\textbf{Second integral:}\\
For $\alpha=0$ it explodes.\\
For $\alpha\neq 0$:
\begin{align}
\int_{\log(K)}^{\infty}e^{-\alpha x}K dx = \lim_{n\to \infty}[-\frac{e^{-\alpha x}}{\alpha}]_{\log(K)}^n<\infty \Leftrightarrow \alpha>0
\end{align}
Combining these results,
\begin{equation}z= \beta+i \alpha \in S_f \Leftrightarrow \alpha \in (0,1)
\end{equation}
Now let's calculate the generalised Fourier transform when it exists.
Let $\alpha \in (0,1)$:
\begin{align}
    \hat{w} (\beta+i \alpha)&=\int_{-\infty}^{\infty}e^{i(\beta+i \alpha)x}w(x) dx\\
    &=\int_{-\infty}^{\log(K)}e^{i(\beta+i \alpha)x}e^x dx + \int_{\log(K)}^{\infty}e^{i(\beta+i \alpha)x}K dx\\
    &=\int_{-\infty}^{\log(K)}e^{i(\beta+i (\alpha-1))x} dx + \int_{\log(K)}^{\infty}e^{i(\beta+i \alpha)x}K dx\\
    &= \lim_{n\to \infty}[\frac{e^{i(\beta+i (\alpha-1))x}}{i(\beta+i (\alpha-1))}]^{\log(K)}_{-n}+K \lim_{n\to \infty}[\frac{e^{i(\beta+i \alpha)x}}{i(\beta+i \alpha)}]^{n}_{\log(K)}\\
    &= \frac{e^{i(\beta+i (\alpha-1))\log(K)}}{i(\beta+i (\alpha-1))} -K \frac{e^{i(\beta+i \alpha)\log(K)}}{i(\beta+i \alpha)}\\
    &= K^{ 1+i(\beta+i \alpha)}(\frac{1}{1+i(\beta+i \alpha)} -\frac{1}{i(\beta+i \alpha)})\\
    &=\frac{K^{1+iz}}{z^2-iz}
\end{align}
where a dominated convergence argument is used on the step from line 17 to line 18, and we swap to $z=\beta+i\alpha$ on the last step.

\subsection{Existence Strip of the CF}

Consider now a storable asset $S_T$ under the risk-neutral measure $\mathbb{Q}$ in a world with constant short rate $r>0$ and continuously paid dividend $q\geq0$. Let
$X_T=\log(S_T)$, a consequence of the risk-neutral measure construction is that $\mathbb{E}[S_T]=\exp((r-q)T)S_0$.
Before moving on we want to show that the CF for basically any type of model is well defined on $\Im(u)\in(0,1)$.\\
In the following everything is happening under $\mathbb{Q}$.\\
Let $u=\beta + i \alpha$.

for $\alpha =0$
\begin{align}
\mathbb{E}[|\exp(-iuX_T)|]=1<\infty
\end{align}

for $\alpha = 1$
\begin{align}
\mathbb{E}[|\exp(-iuX_T)|]&=\mathbb{E}[|\exp(-i(\beta+i\alpha)X_T)|]\\
&=\mathbb{E}[\exp(\alpha X_T)]\\
&=\mathbb{E}[\exp(X_T)]\\
&=\mathbb{E}[S_T]=\exp((r-q)T)S_0<\infty
\end{align}
for $\alpha\in (0,1)$, notice that $x\to x^\alpha$ is concave, then by Jensen's inequality
\begin{align}
\mathbb{E}[\exp(\alpha X_T)]&=\mathbb{E}[(\exp(X_T))^\alpha ]\\
\leq (\mathbb{E}[\exp(X_T)])^\alpha &= S_0^\alpha  \exp(\alpha (r-q)T)<\infty\\
&\Rightarrow \{z\in \mathbb{C} : \operatorname{Im}(z)\in[0,1]\}
\end{align}



\subsection{A General Formula for a Covered-Call Price Using Lewis}
For log-asset dynamics $X_T=\log(S_T)$ under the risk-neutral measure $\mathbb{Q}$,
now assume that
$\{\alpha +\frac{i}{2} : \alpha\in\mathbb{R}\}\subseteq \mathcal{S}_{-\psi} \cap \mathcal{S}_w$. Choosing to integrate along the line $\{\alpha +\frac{i}{2} : \alpha\in\mathbb{R}\}$ we get a simple formula for the covered call at $t=0$.

\begin{theorem}[Covered-Call price]
Consider the covered call payoff function $w(x)=\min(e^x,K)$ and log-asset dynamics $X_T=\log(S_T)$ under $\mathbb{Q}$. Then the covered call price at $t=0$ is:
\begin{equation}Z =\frac{e^{\frac{k}{2}-rT}}{\pi}\int_{0}^{\infty}\frac{\operatorname{Re}[\psi(u-\frac{i}{2})e^{-i k u}]}{u^2+\frac{1}{4}}du
\end{equation}
where $k=\log(K)$ and $\psi$ is the extended characteristic function of $X_T$
\end{theorem}
\begin{proof}
By assumption $\{\alpha +\frac{i}{2} : \alpha\in\mathbb{R}\}\subseteq \mathcal{S}_{-\psi} \cap \mathcal{S}_w$.
We choose to integrate across the line $\{\alpha +\frac{i}{2} : \alpha\in\mathbb{R}\}$, i.e.\ pick $\alpha =1/2$ in the Lewis formula. Taking the expression for $\hat{w}$ that we found:
\begin{align}
    \mathbb{E}[\min(e^{X_T},K)]&=\frac{1}{2\pi}\int_{-\infty+i\alpha}^{\infty+i\alpha}\psi(-z)\hat{w}(z) dz\quad\\
    &=\frac{1}{2\pi}\int_{-\infty+i\alpha}^{\infty+i\alpha}\psi(-z)\frac{K^{1+iz}}{z^2-iz}dz \quad\\
    &=\frac{1}{2\pi}\int_{-\infty+i\alpha}^{\infty+i\alpha}\mathbb{E}[e^{-izX_T}]\frac{K^{1+iz}}{z^2-iz}dz\\
    &=\frac{K}{2\pi}\int_{-\infty+i\alpha}^{\infty+i\alpha}\mathbb{E}[e^{-izX_T}]\frac{K^{iz}}{z^2-iz}dz
\end{align}
now we change variable to $u=z-i\alpha$, note:
\begin{align}
(u+\frac{i}{2})^2-i(u+\frac{i}{2})&=u^2+iu-\frac{1}{4} -iu+\frac{1}{2}\\
&=u^2+\frac{1}{4}\\
\frac{K}{2\pi}\int_{-\infty+i\alpha}^{\infty+i\alpha}\mathbb{E}[e^{-izX_T}]\frac{K^{iz}}{z^2-iz}dz
&=\frac{K}{2\pi}\int_{-\infty}^{\infty}\mathbb{E}[e^{-iuX_T}e^{\frac{X_T}{2}}]\frac{K^{iu}K^{-\frac{1}{2}}}{u^2+\frac{1}{4}}du\\
&=\frac{K^{\frac{1}{2}}}{2\pi}\int_{-\infty}^{\infty}\mathbb{E}[e^{-iuX_T}e^{\frac{X_T}{2}}]\frac{K^{iu}}{u^2+\frac{1}{4}}du
\end{align}
setting $k = \log(K)$
\begin{align}
    &=\frac{e^{\frac{k}{2}}}{2\pi}\int_{-\infty}^{\infty}\psi(-u-\frac{i}{2})\frac{e^{i k u}}{u^2+\frac{1}{4}}du
\end{align}
now we notice:
\begin{align}
    \overline{\psi(-u-\frac{i}{2})}&=\overline{\mathbb{E}[e^{-iuX_T}e^{\frac{X_T}{2}}]}\\
    &=\mathbb{E}[\overline{e^{-iuX_T}e^{\frac{X_T}{2}}}]\\
    &=\mathbb{E}[e^{iuX_T}e^{\frac{X_T}{2}}]\\
    &=\psi(u-\frac{i}{2})\\
    \overline{e^{iku}}&=e^{-iku}\\
    \overline{\frac{1}{u^2+\frac{1}{4}}}&=\frac{1}{(-u)^2+\frac{1}{4}}
\end{align}
combining with the result,
\begin{align}
f\in L^1 \overline{f(x)}=f(-x)\Rightarrow \int_{-\infty}^{\infty}f(x)dx = 2\int_{0}^{\infty}\operatorname{Re}[f(x)]dx=2\int_{0}^{\infty}\operatorname{Re}[f(-x)]dx
\end{align}
we then have
\begin{align}
    \frac{e^{\frac{k}{2}}}{2\pi}\int_{-\infty}^{\infty}\psi(-u-\frac{i}{2})\frac{e^{i k u}}{u^2+\frac{1}{4}}du
    &=\frac{e^{\frac{k}{2}}}{\pi}\int_{0}^{\infty}\frac{\operatorname{Re}[\psi(u-\frac{i}{2})e^{-i k u}]}{u^2+\frac{1}{4}}du
\end{align}
Discounting then gives
\begin{equation}
Z = e^{-rT}\mathbb{E}[\min(e^{X_T},K)] = \frac{e^{\frac{k}{2}-rT}}{\pi}\int_{0}^{\infty}\frac{\operatorname{Re}[\psi(u-\frac{i}{2})e^{-i k u}]}{u^2+\frac{1}{4}}du
\end{equation}
\end{proof}
\begin{remark}
We can now price the put and call by replication:
\begin{align}
C &= e^{-rT}F_T-Z = e^{-rT}F_T-\frac{e^{\frac{k}{2}-rT}}{\pi}\int_{0}^{\infty}\frac{\operatorname{Re}[\psi(u-\frac{i}{2})e^{-i k u}]}{u^2+\frac{1}{4}}du\\
P &= Ke^{-rT}-\frac{e^{\frac{k}{2}-rT}}{\pi}\int_{0}^{\infty}\frac{\operatorname{Re}[\psi(u-\frac{i}{2})e^{-i k u}]}{u^2+\frac{1}{4}}du
\end{align}
\end{remark}
